% !TeX root = ../papers/nc-p3.tex
\section{二次錐の三次・五次結果式と中間座標の排除}
\label{sec:quadric-cone-middle}
\begin{theorem}\label{thm:general-quadric-cone-exclusion}
二次錐 $Q$ の $\operatorname{Gr}(4,H^0(Q,\mathcal O_Q(2)))$ のある非空稠密開集合では、
二十次二重曲線の射影正規形をもつ被約有限標準正則窓は、どの中間座標を取っても存在しない。
\end{theorem}
\begin{proof}[錐の中間座標]
$R=k[s,t,w]$, $\deg(s,t,w)=(1,1,2)$ とし、$S_j=k[s,t]_j$ とする。
$V\subset R_4$、核平面は $y=a+bw\in R_2$ で表す。
一般の $V$ は全母線で全二次式を与え、頂点の入力繊維は被約なその一点である。
前者は固定した母線で全射の失敗が余次元二、後者は頂点と他の一点の分離失敗が固定二点で余次元三であることによる
（完全な二次系は非常に豊富）。さらに頂点像での相対余接空間は一般に零である。
実際頂点で消える三座標の一次部は $s^2/w,st/w,t^2/w$ の一般の三一次結合である。

母線に沿う評価は底上の全射列
\[
 0\longrightarrow\mathcal O(-6)\longrightarrow V\otimes\mathcal O
 \longrightarrow\mathcal O(4)\oplus\mathcal O(2)\oplus\mathcal O\longrightarrow0
\]
を与える。最後の三成分は $1,w,w^2$ の係数である。$1,0,-1$ 捻りの接続写像と双対性により
$A^*=(R_5/VS_1)^*=S_3$, $(R_4/V)^*=S_4$, $R_3^*=S_5$。
この順に $s,t$ 倍の転置は二項式の積である。$w$ 倍の転置を $D:S_3\to S_5$ とおくと六個の中間辺の式は
\[
       m(yR_3)=0\quad\Longleftrightarrow\quad a m+b Dm=0 .
\]
特に $b\ne0$ であり、印の条件は三次式 $m$ が $Dm$ を割ることである。
もう一方の辺を
$R_{xm}=\ker(V\otimes A\to R_9/((\operatorname{Sym}^2V)S_1))$ とする。
この核が六次元で両辺の両射影が有限である開条件を課す。

実際の二十印で、長さ三繊維 $Y$ は頂点を避ける。各母線上で入力は埋込みだから、
底への各繊維長は一以下で、有限射 $Y\to\mathbb P^1$ は閉埋込みである
（各点の局所環で中山の補題を適用）。$R_5$ は $Y$ の全切断を与える。
実際 $R_4$ は三次元射影空間の二次式の制限で長さ三を分離し、
$Y$ 上で消えない $S_1$ の元を掛けられる。
$VS_1$ の値は二次元なのでその零化汎関数 $m$ が定まり、繊維で消える $y$ と中間辺を満たす。
$R_{xm}$ の関係も繊維上で $V$ の値を単元倍の $x$ に置き換えて汎関数を施せば成立する。
従ってこの二つの有限完全交叉グラフはそれぞれ二十組で尽きる。
特に $x,m$ グラフの逆繊維も線形であり $m$ は相異なる。以下この配置を $M$ と書く。
\end{proof}

\begin{proof}[三次・五次の割線と別の平方根]
$m(z)=(Dm)(z)=0$ は各 $z\in\mathbb P^1$ で $m$ の二一次式である。
この二行を $B(z)$ とし、二点因子の平面に
$P=\det(B(z);B(z'))/(z-z')^2$ を降ろす（分母は斉次には対角式の平方）。
二行の次数の和は八であり、商は対称二次数 $(6,6)$、従って平面六次である。
$P$ が滑らかで対角二次曲線との交点は単純、全ての印の三次式が平方自由、
$M$ の三重八次式の空間は一次元と仮定する。後で開条件を検証する。

定理\ref{thm:generic-quadric-all-middles-exclusion}の割線の議論がこの二行にも適用できる。
実際ある $B(z)$ の階数が落ちれば $P$ は直線成分を持ち、相異なる二母線が一致すれば $P$ は特異になる。
母線束からの射は有限である（各 $m$ の逆像は非零三次式の根に限られる）。
三母線に乗る点は $m\mid Dm$ より有限。従って像は整な結果式八次曲面
$D_0=V(\operatorname{Res}(m,Dm))$ で、この母線束は正規化である。
結果式は一般の単純な一共通根で一重である。これは母線束の超平面次数が
基点のないPlücker八次式の次数八であることからも従う。
各印で三つの単純な共通根があるため結果式は三重に消える
（局所エタールに三根を分離して積に書く）。ゆえに別の非恒等根 $N$ の行列式は
定理\ref{thm:opposite-root-triple-octics}の被約配置での階数排除によりこの結果式である。
ここで $M$ の三次評価同型と四点直線の排除は有限辺から従う。

$P$ は唯一の特異曲線 $C$ の正規化である。相異なる母線対は結合行列の階数三の小行列で
交点を定め、三本が共点となる有限集合の外ではその交点から母線対が一意に戻る。
$P$ の滑らかさから一般二枝は横断し、対角制限の単純性から非浸入点は有限となる
（定理\ref{thm:generic-quadric-all-middles-exclusion}の微分判定）。
一般平面八次切断は有理で節点のみを持ち、従って $\deg C=21$。
各印に対応する三根の三つの異なる二点因子は $P$ の相異なる三点である。

$N$ が $C$ の一般点で階数三と仮定する。$P$ へ引き戻した六次余因子は次数 $126$ の束の切断で、
二十印の上の六十点それぞれで二重に消える。共通固定部を引いた系の次数は従って六以下である。
この射は核点の二次Veronese像に入り、核写像の次数（超平面束の次数）は三以下になる。
一方滑らかな平面六次上の次数三以下の有効因子 $E$ は動かない。
実際標準束は平面三次式であり、長さ三以下の全ての $E$ の値を分離するので
Riemann--Rochで $h^0(E)=1$。ゆえに核点は一定である。
すると $N$ の固定ベクトルとの積が $C$ 上で消え、その非零二次成分が二十印を通ることになり矛盾する。

従って $C$ の一般点で階数二でなければならない。節点の二次初項から余因子曲線はそこで被約。
定理\ref{thm:intrinsic-secant-scroll-root}の正次元の場合の証明より、一般平面で核次数は三、
核像は有理正規三次であり、一般核繊維は既存の一意な母線族である。
同じEuler数の比較から $\operatorname{coker}N=\nu_*\pi^*\mathcal O(3)$
（$\pi$ は母線の底、$\nu$ は正規化写像）。
ここで四生成元は底の完全三次系そのものである。しかし $m(z)=0$ は全母線束上でこれらの
四切断の非零な入力一次係数関係である。四行の二次式による極小表示には一次関係がない。よって非恒等根は存在しない。
\end{proof}

\begin{proof}[折り返しの交代形式]
実際の二つの中間配置は、有限辺の完全跡の直交性から、ともに $M$ の平方配置の根となる。
上で示した一意性より両者とも $M$ である。有限完全交叉の双一次式は六個の定義式で尽きるので
中央二辺は $R_{my}$ とその転置である。命題\ref{prop:finite-flip-obstructions}により頂テンソルは
二つの $m$ で対称である。この折り返しの対称三次交叉が四以下で、四ならそこから得る頂テンソルが
反対階数十を持つ条件を課す。これは後述の標本の近傍で、正則候補すべてに対して成立する。
実際線形方程式の独立な二十行で二十四未知数のうち二十を消去し、残る四パラメータからの反対切断の
十行小行列の非消滅を課せばよい。正則なら四次元全体が解になる。

この辺は $D$ によらず $m$ 側の平方根接空間に階数六の方向を持つ。
実際 $b$ に対応する文字を割ると関係は $S_5^*\to S_3^*\otimes S_2^*$、
二項乗法の転置となる。商の次元は六で、その上に非退化な交代形式で
$S_3^*$ 側に交代、$S_2^*$ 側に対称なものが存在する。
これは固定された整数行列で確認できる。$C_{(i,j),l}=\delta_{i+j,l}$ とし
$B$ の列を $\ker C^t$ の基底とする。交代六行列 $J$ について
$(BJB^t)_{ui,vj}+(BJB^t)_{vi,uj}=0$ を全 $u\le v,i<j$ で解くと
解は一次元で非零解の階数は六である（下記検算）。
元の十次元商への引戻しも階数六で、定理\ref{thm:finite-square-tangent-instantons}の接条件を満たす。
定理\ref{thm:single-edge-instanton-extension}により反対階数十の正則折り返しなら
全ての非零方向は階数十であり、矛盾する。

\path{computations/verify_quadric_cone_middle.py} は一つの有理四空間から両辺を構成する。
両向きの平面上の十五最大小行列式、滑らかな割線六次と十二個の単純対角交点、
二十点の被約性と全三次判別式の可逆性、八次三重評価階数164を素数1009で検算する。
八次結果式の存在が最後の階数の反対向きの評価を与える。
対称折り返しの四次元核と反対階数十、固定商の交代形式は有理数上で調べる。
いずれも指定された開条件を満たし、母線・頂点の一般条件との共通開集合もGrassmann多様体の既約性から非空である。
\end{proof}
